\section{The reconstruction problem}
\label{sec:introduction}

Birational anabelian geometry asks how much of a field can be read from its
Galois groups.  Over algebraically closed constants, the arithmetic Galois
action disappears.  Bogomolov's program proposes that a particularly small
quotient still remembers every function field of dimension at least two:
the quotient in which commutators are central.  We prove the Isom form of
this assertion, including surfaces and the prime $2$.

\subsection{The exact statement}
Fix a prime $\ell$.  A function field $F/\kappa$ in this paper is a finitely
generated extension of an algebraically closed field $\kappa$.
For $\operatorname{char}F\ne\ell$, let $G_F^{(\ell)}$ be the maximal
pro-$\ell$ quotient of the absolute Galois group, and set
\[
 \Pi_F^a=G_F^{(\ell)}/[G_F^{(\ell)},G_F^{(\ell)}],\qquad
 \Pi_F^c=G_F^{(\ell)}/[[G_F^{(\ell)},G_F^{(\ell)}],G_F^{(\ell)}].
\]
All commutator subgroups are closed.  The kernel
$\Delta_F=\ker(\Pi_F^c\to\Pi_F^a)$ is central.  Our commutator convention is
$[x,y]=x^{-1}y^{-1}xy$.  Lifting elements of $\Pi_F^a$ to $\Pi_F^c$ therefore
defines a continuous alternating $\Zell$-bilinear map
\[
 [\ ,\ ]_F:\Pi_F^a\times\Pi_F^a\longrightarrow\Delta_F.
\]

An isomorphism $\varphi:\Pi_L^a\to\Pi_K^a$ is
\emph{bracket-compatible} if it is a continuous $\Zell$-module isomorphism
and there is a continuous $\Zell$-module isomorphism
$\psi:\Delta_L\to\Delta_K$ such that
\[
 \psi([\sigma,\tau]_L)=[\varphi\sigma,\varphi\tau]_K
 \quad(\sigma,\tau\in\Pi_L^a).
\]
Write $\Isom^c(\Pi_L^a,\Pi_K^a)$ for this set.  The units $\Zell^\times$
act by multiplication on $\varphi$; the corresponding map on the central
kernel can be multiplied by the square of the same unit.

Write $F^i$ for the perfect closure of $F$; in characteristic zero this is
$F$ itself.  Let $\Isom^i(K,L)$ consist of the field isomorphisms
$\alpha:K^i\to L^i$ satisfying $\alpha(k)=l$.
If the common characteristic is $p>0$, put
\[
 \Isom_F^i(K,L)=\Isom^i(K,L)/
 \bigl(\alpha\sim\Frob_L^n\circ\alpha,\ n\in\mathbb Z\bigr).
\]
Frobenius is an automorphism of a perfect field, so negative integers are
allowed.  In characteristic zero take no quotient.  If the
characteristics differ, set $\Isom_F^i(K,L)=\varnothing$.

Field isomorphisms induce contravariant Galois isomorphisms.  Passing to the
above quotients gives the canonical map in the following theorem.

\begin{theorem}[Bogomolov--Pop reconstruction]\label{thm:main}
Let $k$ and $l$ be arbitrary algebraically closed fields of characteristics
different from $\ell$, and let $K/k$ and $L/l$ be function fields with
$\trdeg(K/k),\trdeg(L/l)\ge2$.  Then
\[
 \Phi_{K,L}:\Isom_F^i(K,L)\longrightarrow
 \Isom^c(\Pi_L^a,\Pi_K^a)/\Zell^\times,
 \qquad [\alpha]\longmapsto[\alpha^*],
\]
is a bijection.
\end{theorem}

Thus the theorem resolves the pro-$\ell$ abelian-by-central
Bogomolov--Pop reconstruction conjecture positively.  Neither equality of
characteristics nor equality of relative transcendence degrees is assumed.
Both are forced by a bracket-compatible isomorphism.  The input contains
no distinguished valuations, inertia groups, or curve quotients.

\subsection{Previous work and the proof}
Bogomolov introduced the abelian-by-central approach to birational
reconstruction \cite{Bogomolov1991}.  The local approach detects
valuations through commuting elements and their associated multiplicative
characters \cite{BT2002,TopazCL2017}.
Bogomolov and Tschinkel proved the surface case over algebraic closures
of finite fields \cite{BT2008}, followed by higher-dimensional
reconstruction \cite{BT2011}; Pop established a corresponding
birational reconstruction theorem \cite{PopBI2012}.
The precise Isom formulation, including its scalar and Frobenius
ambiguities, is recorded in
\cite[Introduction, Conjecture~1]{TopazRat2016}.

Pop's global reconstruction
theorem recovers a function field from its total decomposition graph
together with suitable rational quotients
\cite{PopGraphs2012,PopBI2012}.  His later reconstruction theorem with
divisorial inertia applies in relative transcendence degree greater than
two \cite[Theorem~1.1]{Pop2018}; its additional inertia datum must itself
be recovered to solve the conjecture from the group alone.
The local theory of commuting pairs supplies quasi-divisorial valuations,
which can be nontrivial on the constants
\cite{Pop2018,Topaz2019}.  Pop's minimized decomposition theory also
recognizes certain residual curves over algebraic closures of finite
fields, with Topaz's appendix treating the required residue
characteristic $\ell$ case \cite{PopMin2015}.

Related reconstruction results use additional structure.
Silberstein studies the full absolute Galois group under a
prime-field-closure hypothesis \cite{Silberstein2013}.
Topaz's integral cohomological reconstruction uses a specified
arithmetic Galois action in transcendence degree at least three
\cite[Theorems~A and~B]{Topaz2019}; his Torelli-type theorem uses
mixed Hodge structures and cup products over constants embedded in
$\mathbb C$ \cite{TopazTorelli2025}.
These results concern different input data.  Here only the original
pro-$\ell$ abelian-by-central datum is supplied, and neither the
constant field nor geometric inertia is distinguished in advance.

Our task is to recover the missing geometric information uniformly over
the constants.  Write $\widehat F=\varprojlim_e F^\times/F^{\times\ell^e}$
for the completed multiplicative group and $[f]$ for the class of a
nonzero function.  Choose Kummer identifications, and dualize a
bracket-compatible isomorphism to an isomorphism
$\Theta:\widehat K\to\widehat L$ of completed multiplicative groups.
The central step proves that, for fixed $t\in K\setminus k$, the entire
pencil
\[
 h_a=\Theta[t-a],\qquad a\in k,
\]
has uniformly bounded divisor-support degree on a fixed projective model
of $L$.  This is stronger than finiteness of each individual support.

Bogomolov--Tschinkel already used bounded-genus curves and Hurwitz
estimates to prove finiteness of $\ell$-adic supports
\cite[Section~6, proof of Proposition~6.1]{BT2011}.
Here the bound is uniform over a whole pencil and is obtained by finite
tests over arbitrary constants.  A finite list of Kummer classes
is specialized and restricted to a smooth residual curve of bounded
genus, preserving all relations at the chosen finite level and the
relevant divisor supports.  Five fixed members of the pencil then bound
the degree of a separable map to the projective line by Riemann--Hurwitz.
The bound also controls every additional member, at every sufficiently
high $\ell$-power level.  The use of five anchors works for $\ell=2$.

The pencil detects infinitely many prime divisors, and the uniform
degree bound places them in a family of finite type.  Their incidence
with the model produces a curve subfield
after a finite field extension.  A normal-base-change argument allows us
to use one parameter curve for every Kummer class and every exponent;
no uncountability assumption enters.  Norms and the finite-rank
intersection of distinct completed curve fields descend this subfield
to $L$.  Applying the same argument to $\Theta^{-1}$ gives a bijection
of the completed relatively algebraically closed curve subfields.

These subfields distinguish true divisorial valuations from those that
also value the constants.  They also recover point inertia and genus,
and hence all relatively algebraically closed rational subfields.
Their quotient diagrams supply exactly the input to Pop's global
reconstruction theorem.  The new geometric steps are the uniform
pencil bound and the incidence argument that recovers curve subfields;
the remaining steps use the local valuation theory and Pop's
reconstruction theorem described above.

The proof follows this order.  Section~\ref{sec:preliminaries} fixes the
Kummer topology and the relation between the two formulations of the
group input.  Section~\ref{local:section} supplies the local correspondence.
Sections~\ref{sec:tests} and~\ref{sec:pencil} prove the residual test and
uniform support bound.  Section~\ref{sec:curve-recovery} recovers curve
subfields and rational quotients.  Section~\ref{sec:reconstruction}
checks compatibility on the decomposition graphs before applying Pop's
theorem and completing the Isom statement and its ambiguity checks.
